\documentclass[11pt]{article}
\usepackage{enumitem}
\usepackage{amsthm}
\usepackage{amssymb, amsmath, xcolor}
\usepackage{geometry}
\usepackage[numbers,square,sort&compress]{natbib}
\usepackage{url}
\usepackage[hidelinks]{hyperref}
\usepackage[capitalise,noabbrev]{cleveref}
\usepackage{xspace}
\usepackage{booktabs}
\usepackage{graphicx}

% Draft comment commands
\newcommand{\claude}[1]{\textcolor{red}{\textbf{[Claude:} #1\textbf{]}}}

\newcommand{\Lean}{\textsc{Lean}~4\xspace}
\newcommand{\Mathlib}{\textsc{Mathlib}\xspace}
\newcommand{\A}{\mathfrak{A}}
\newcommand{\norm}[1]{\left\|#1\right\|}
\newcommand{\abs}[1]{\left|#1\right|}
\newcommand{\comm}[2]{[#1,\,#2]}
\DeclareMathOperator{\tr}{tr}

%% --- Hyperlinked Lean references ---
%% \leanref{File}{line}{name} renders as clickable \texttt{name}
%% linking to GitHub source.  \leansrc{File} links to a whole file.
\newcommand{\repourl}{https://github.com/Jue-Xu/Lean-Trotter/blob/main}
\newcommand{\leanref}[3]{\href{\repourl/LieTrotter/#1\#L#2}{\texttt{#3}}}
\newcommand{\leansrc}[1]{\href{\repourl/LieTrotter/#1}{\texttt{#1}}}

\geometry{margin=1in}

%% --- Lean code blocks ---
%% For local builds: xelatex -shell-escape main.tex
%% For arXiv: switch to [frozencache] once the _minted cache is generated.
\usepackage{listings}
\lstdefinelanguage{Lean4}{
  morekeywords={theorem,lemma,def,noncomputable,section,variable,open,import,
    private,where,by,calc,exact,rw,simp,intro,apply,refine,obtain,have,let,set,
    fun,in,if,then,else,match,with,instance,class,structure,namespace,end,
    include,omit,induction,cases},
  sensitive=true,
  morecomment=[l]{--},
  morecomment=[s]{/-}{-/},
  morestring=[b]",
  literate=
    {𝕂}{{$\mathbb{K}$}}1
    {𝔸}{{$\mathfrak{A}$}}1
    {ℝ}{{$\mathbb{R}$}}1
    {ℂ}{{$\mathbb{C}$}}1
    {ℕ}{{$\mathbb{N}$}}1
    {ℚ}{{$\mathbb{Q}$}}1
    {→}{{$\to$}}1
    {←}{{$\leftarrow$}}1
    {↦}{{$\mapsto$}}1
    {≤}{{$\le$}}1
    {≥}{{$\ge$}}1
    {≠}{{$\ne$}}1
    {⁻¹}{{$^{-1}$}}1
    {•}{{$\bullet$}}1
    {∀}{{$\forall$}}1
    {∃}{{$\exists$}}1
    {∈}{{$\in$}}1
    {∑}{{$\sum$}}1
    {∫}{{$\int$}}1
    {∞}{{$\infty$}}1
    {‖}{{$\|$}}1
    {·}{{$\cdot$}}1
}
\lstset{
  language=Lean4,
  basicstyle=\small\ttfamily,
  keywordstyle=\bfseries\color{blue!70!black},
  commentstyle=\itshape\color{green!50!black},
  stringstyle=\color{red!60!black},
  numbers=left,
  numberstyle=\tiny\color{gray},
  numbersep=6pt,
  frame=lines,
  framesep=2mm,
  breaklines=true,
  tabsize=2,
  showstringspaces=false
}
\lstnewenvironment{leancode}{\lstset{language=Lean4}}{}

\title{Formalizing the Lie--Trotter Product Formula
and Commutator-Scaling Error Bounds in Lean~4}
% \author{Jue Xu}

\newtheorem{theorem}{Theorem}[section]
\newtheorem{lemma}[theorem]{Lemma}
\newtheorem{proposition}[theorem]{Proposition}
\newtheorem{corollary}[theorem]{Corollary}

\theoremstyle{definition}
\newtheorem{definition}[theorem]{Definition}

\theoremstyle{remark}
\newtheorem*{remark}{Remark}

\begin{document}
\maketitle

\begin{abstract}
We present a machine-checked proof of the Lie--Trotter product formula
$e^{A+B} = \lim_{n\to\infty} (e^{A/n}\,e^{B/n})^n$
for elements $A$, $B$ in a complete normed algebra, formalized in \Lean with the
\Mathlib library.
Beyond the basic convergence theorem with an explicit $O(1/n)$ error bound,
our development includes the symmetric Strang splitting with $O(1/n^2)$ convergence,
multi-operator generalizations, a fourth-order Suzuki integrator,
and---most significantly---commutator-scaling error bounds matching those of
Childs et al.\ (2021), where the product $\norm{A}\norm{B}$ is replaced by
the commutator norm $\norm{\comm{B}{A}}$.
The commutator-scaling proofs introduce a Duhamel/FTC-based technique that
avoids ODE uniqueness theory entirely.
The formalization comprises over 4\,300 lines of Lean~4 with zero
\texttt{sorry} axioms, building against current \Mathlib and filling several
gaps in the library's treatment of the matrix exponential.
\end{abstract}

\section{Introduction}
\label{sec:intro}

Product formulas are among the most widely used tools in quantum simulation.
The Lie--Trotter formula
\begin{equation}\label{eq:lie-trotter}
  e^{A+B} = \lim_{n\to\infty} \bigl(e^{A/n}\,e^{B/n}\bigr)^n
\end{equation}
decomposes the exponential of a sum into a product of simpler exponentials,
enabling the simulation of composite quantum systems by sequentially
simulating each term of the Hamiltonian~\cite{trotter1959product, lloydUniversalQuantumSimulators1996}.
Higher-order variants---the symmetric Strang splitting~\cite{strang1968construction}
and Suzuki's recursive hierarchy~\cite{suzukiGeneralTheoryFractal1991}---achieve
faster convergence by exploiting symmetry and commutator cancellation.
Rigorous error bounds for these formulas are critical for resource estimation
in quantum algorithms~\cite{childsTheoryTrotterError2021}.

The quality of these error bounds directly affects the efficiency of quantum
algorithms.  The standard bound scales with the product $\norm{A}\norm{B}$,
which treats $A$ and $B$ as unrelated.
Childs et al.~\cite{childsTheoryTrotterError2021} initiated a program of
\emph{commutator-scaling} bounds, showing that the leading error is controlled
by the commutator norm $\norm{\comm{B}{A}}$ rather than the product
$\norm{A}\norm{B}$.
For spatially local Hamiltonians $H = \sum_j h_j$ typical of many-body physics,
most pairs of terms commute, so $\sum_{j<k} \norm{\comm{h_j}{h_k}}$ can be
polynomially smaller than $\sum_{j<k} \norm{h_j}\norm{h_k}$.
These tighter bounds translate directly into fewer Trotter steps---and hence
fewer quantum gates---for a given simulation accuracy.

In this work, we present a complete formalization of the Lie--Trotter product
formula and its higher-order extensions in \Lean, an interactive theorem prover
with a small trusted kernel~\cite{demoura2021lean4}.
The \Mathlib library provides a large corpus of formalized
mathematics~\cite{mathlib2020}, including the normed algebra exponential,
upon which our development builds.

Proofs of product formula error bounds involve long chains of norm estimates
in non-commutative algebras, where each step interacts with the previous through
submultiplicativity, the triangle inequality, and exponential series manipulation.
A single bookkeeping error---a missing factor, a flipped sign in a
commutator---can silently propagate and invalidate the final bound.
Formal verification forces every intermediate step to be explicit and
machine-checked, making it a natural fit for this domain.

Our formalization proves the following hierarchy of results:
\begin{enumerate}
  \item \textbf{First-order Lie--Trotter} (\cref{sec:assembly}):
    $\norm{(e^{A/n}\,e^{B/n})^n - e^{A+B}} \le C/n$
    with explicit constant $C = 2\norm{A}\norm{B}e^{\norm{A}+\norm{B}} + 1$.

  \item \textbf{Strang splitting} (\cref{sec:strang}):
    $\norm{(e^{A/2n}\,e^{B/n}\,e^{A/2n})^n - e^{A+B}} \le C'/n^2$.

  \item \textbf{Multi-operator extensions}:
    $(e^{A_1/n} \cdots e^{A_m/n})^n \to e^{A_1+\cdots+A_m}$
    and the palindromic (Strang) variant, both with explicit error bounds.

  \item \textbf{Commutator-scaling bounds} (\cref{sec:commutator}):
    \begin{align}
      \norm{e^{tB}\,e^{tA} - e^{t(A+B)}}
        &\le \tfrac{1}{2}\norm{\comm{B}{A}}\,t^2\,e^{t(\norm{A}+3\norm{B})},
        \label{eq:comm-first} \\
      \norm{S_2(t) - e^{tH}}
        &\le \bigl(\tfrac{1}{12}\norm{\comm{B}{\comm{B}{A}}}
          + \tfrac{1}{24}\norm{\comm{A}{\comm{A}{B}}}\bigr)\,t^3,
        \label{eq:comm-strang}
    \end{align}
    matching the commutator-scaling program of Childs et al.~\cite{childsTheoryTrotterError2021}.

  \item \textbf{Tighter Strang bound via norm-of-difference}: a strictly
    sharper second-order bound
    $\norm{S_2(t)-e^{tH}} \le \tfrac{\norm{D}}{6}t^3 + \tfrac{T}{4}t^4$
    where $D = \comm{B}{\comm{B}{A/2}} - \comm{A/2}{\comm{A/2}{B}}$ is an
    effective double commutator capturing signed cancellations that the
    triangle-inequality form of \eqref{eq:comm-strang} cannot see.

  \item \textbf{Fourth-order Suzuki $S_4$ with CAS-assisted BCH bounds}
    (\cref{sec:suzuki4}):
    four formalized order-$t^5$ bounds for
    $\norm{S_4(t) - e^{t(A+B)}}$, culminating in a provably
    strictly tighter (9$\times$--64$\times$ per coefficient)
    replacement of Childs's~2021 heuristic prefactors, together with an
    existential-$\delta$ finite-$t$ bound obtained by extending the CAS
    computation to order~$\tau^7$.

  \item \textbf{Fourth-order convergence of $S_4$} (\cref{sec:s4-convergence}):
    compounding the single-step $O(t^5)$ bound over $n$ steps gives
    $\norm{(S_4(t/n))^n - e^{t(A+B)}} \le C/n^4$ for $n$ beyond an explicit
    threshold, and hence $(S_4(t/n))^n \to e^{t(A+B)}$. Together with items~1
    and~2 this completes a verified convergence hierarchy
    $O(1/n) \to O(1/n^2) \to O(1/n^4)$.
\end{enumerate}

The entire development comprises 32 Lean source files totaling over 12\,100 lines,
with zero \texttt{sorry} axioms and zero own theorem-level BCH-interface
axioms on the Lean-Trotter side. The three former BCH-interface axioms
(\texttt{bch\_w4Deriv\_quintic\_level2},
\texttt{bch\_w4Deriv\_level3\_tight},
\texttt{bch\_uniform\_integrated})
are now theorems composing Lean-BCH bridge corollaries with
exp-Lipschitz lifts. As of Lean-BCH revision \texttt{d455ff0}
(2026-05-19), the upstream B1.c quintic axiom has been discharged, so
all $\tau^5$ Suzuki~$S_4$ headlines---together with the standard
Lie--Trotter, Strang, and commutator-scaling theorems---depend only
on the standard foundational axioms \texttt{[propext, Classical.choice,
Quot.sound]}. In particular, axiom inspection on
\texttt{norm\_suzuki4\_childs\_form\_via\_level3} (Level~1 Childs
reproduction), \texttt{norm\_suzuki4\_level3\_bch} (Level~3 tight
$\gamma_i$ prefactors), \texttt{norm\_suzuki4\_level2\_bch} (Level~2
unit-coefficient $\tau^5$ bound), the abstract $S_4$ $O(t^5)$ identity
(\texttt{bch\_iteratedDeriv\_s4Func\_order4}), and \texttt{lie\_trotter}
returns the three foundational axioms only. Two transitive private
axioms remain inside the companion
\href{https://github.com/Jue-Xu/Lean-BCH}{Lean-BCH} project---both
septic stepping stones
(\texttt{symmetric\_bch\_septic\_sub\_poly\_axiom},
\texttt{norm\_septic\_match\_residual\_le\_axiom}). They gate only
the optional Level~4 finite-$t$ uniform refinement with an explicit
$\tau^7$ correction, not the core tight 4th-order bound. All other
Lean-BCH content is fully sorry-free; the three-factor symmetric BCH
cubic identities, the two-factor palindromic quintic remainder, and
the five-factor palindromic quintic identification are theorems
specialized to $\mathbb{K} = \mathbb{R}$ and imported directly.
Our \Lean formalization is available at
\href{https://github.com/Jue-Xu/Lean-Trotter}{\texttt{github.com/Jue-Xu/Lean-Trotter}}.
Our main contributions are:
\begin{itemize}
  \item The \emph{first formalization} of the Lie--Trotter product formula in any
    proof assistant, including its higher-order and multi-operator extensions.

  \item A \emph{Duhamel/FTC-based proof technique} for commutator-scaling bounds
    that avoids ODE uniqueness theory (Gronwall's inequality, Picard iteration),
    requiring only the fundamental theorem of calculus for Bochner integrals.

  \item Machine-checked \emph{commutator-scaling error bounds} at first and
    second order, with double-commutator coefficients matching the
    theoretical literature---and a new tighter-Strang refinement
    via norm-of-difference (\cref{apd:tighter}).

  \item \emph{CAS-assisted BCH bounds for Suzuki $S_4$}: four machine-checked
    bounds---including a rigorously tighter replacement for the
    Childs (2021) prefactors and an existential-$\delta$ finite-$t$ bound
    via an order-$\tau^7$ correction---with all CAS output verified by Lean
    numerical sanity checks.

  \item Identification and filling of several \emph{gaps in \Mathlib}, notably
    the inequality $\norm{e^a} \le e^{\norm{a}}$ for general Banach algebras.
\end{itemize}

To our knowledge, no prior formalization of the Lie--Trotter product formula
exists in any proof assistant.
\Mathlib contains the identity $e^{a+b} = e^a\,e^b$ for commuting elements
(\texttt{exp\_add\_of\_commute}), but the non-commuting convergence
theorem---which is a limit statement, not an algebraic identity---is new.

The remainder of this paper is organized as follows.
\Cref{sec:background} reviews the mathematical background on product formulas
and the Lean/\Mathlib ecosystem.
\Cref{sec:proof} presents the proof architecture, including the modular
decomposition and the Duhamel technique for commutator scaling.
\Cref{sec:formalization} discusses the formalization in \Lean, with
representative code excerpts and design decisions.
\Cref{sec:discussion} reflects on the development process, including
AI-assisted formalization and lessons learned.
\Cref{sec:conclusion} concludes with future directions.

\section{Mathematical Background}
\label{sec:background}

\subsection{The Lie--Trotter product formula}
\label{sec:bg-trotter}

Let $\A$ be a unital Banach algebra over $\mathbb{K} \in \{\mathbb{R}, \mathbb{C}\}$
with norm $\norm{\cdot}$ satisfying $\norm{1} = 1$ and $\norm{ab} \le \norm{a}\norm{b}$.
The exponential
\[
  e^a = \sum_{k=0}^{\infty} \frac{a^k}{k!}
\]
converges absolutely for every $a \in \A$.

Trotter~\cite{trotter1959product} proved that for bounded operators $A, B$,
\[
  e^{A+B} = \lim_{n\to\infty} \bigl(e^{A/n}\,e^{B/n}\bigr)^n.
\]
The standard proof proceeds in two steps:
\begin{enumerate}
  \item \emph{Telescoping}: writing $M = \max\{\norm{X},\norm{Y}\}$,
    the identity
    $X^n - Y^n = \sum_{k=0}^{n-1} X^k(X-Y)Y^{n-1-k}$ gives
    $\norm{X^n - Y^n} \le n\,\norm{X-Y}\,M^{n-1}$.

  \item \emph{Step error}: the single-step approximation satisfies
    $\norm{e^a\,e^b - e^{a+b}} = O(\norm{a}\,\norm{b})$.
\end{enumerate}
Setting $X = e^{A/n}\,e^{B/n}$ and $Y = e^{(A+B)/n}$, the $O(1/n^2)$ step
error is amplified by the factor $n$ from telescoping, yielding $O(1/n)$
overall convergence.

\subsection{Higher-order splitting methods}
\label{sec:bg-splitting}

The symmetric \emph{Strang splitting}~\cite{strang1968construction} uses
\[
  S_2(t) = e^{tA/2}\,e^{tB}\,e^{tA/2}
\]
to achieve $O(t^3)$ local error and $O(1/n^2)$ global convergence.
Suzuki~\cite{suzukiGeneralTheoryFractal1991} showed that higher-order integrators can be
constructed recursively: the fourth-order formula
\[
  S_4(t) = S_2(pt)^2\,S_2\bigl((1-4p)t\bigr)\,S_2(pt)^2
  \qquad \text{with } p = \frac{1}{4 - 4^{1/3}}
\]
achieves $O(t^5)$ local error.

Both formulas generalize to multiple operators $A_1, \ldots, A_m$.
The first-order product $(e^{A_1/n} \cdots e^{A_m/n})^n$ converges at rate
$O(1/n)$, while the palindromic product
\[
  S_2^{(m)}(t) = e^{tA_1/2}\,e^{tA_2/2}\cdots e^{tA_m}\cdots e^{tA_2/2}\,e^{tA_1/2}
\]
achieves $O(1/n^2)$.

\subsection{Commutator-scaling error bounds}
\label{sec:bg-commutator}

The standard step error $\norm{e^a\,e^b - e^{a+b}} \le 2\norm{a}\norm{b}e^{\norm{a}+\norm{b}}$
treats $A$ and $B$ as unrelated. When $A$ and $B$ nearly commute---as is
typical for spatially local Hamiltonians---the error is much smaller.

Childs et al.~\cite{childsTheoryTrotterError2021} showed that the leading error term is
controlled by the commutator $\comm{B}{A} = BA - AB$ rather than the product
$\norm{A}\norm{B}$. At first order,
\begin{equation}\label{eq:comm-scaling-bg}
  \norm{e^{tB}\,e^{tA} - e^{t(A+B)}}
    \le \frac{\norm{\comm{B}{A}}}{2}\,t^2\,e^{t(\norm{A}+3\norm{B})}.
\end{equation}
For the Strang splitting, the leading commutator $\comm{B}{A}$ cancels by
symmetry, and the error is controlled by double commutators:
\begin{equation}\label{eq:strang-comm-bg}
  \norm{S_2(t) - e^{tH}}
    \le \Bigl(\frac{\norm{\comm{B}{\comm{B}{A}}}}{12}
      + \frac{\norm{\comm{A}{\comm{A}{B}}}}{24}\Bigr)\,t^3.
\end{equation}

These commutator-scaling bounds are physically significant: for a lattice
Hamiltonian $H = \sum_j h_j$ where each $h_j$ acts on a few neighboring sites,
$\norm{\comm{h_j}{h_k}} = 0$ whenever the supports of $h_j$ and $h_k$ are
disjoint, so the commutator-based bound can be polynomially tighter than the
naive product bound~\cite{childsTheoryTrotterError2021}.

\subsection{Lean~4 and Mathlib}
\label{sec:bg-lean}

\Lean is an interactive theorem prover and programming language based on
dependent type theory~\cite{demoura2021lean4}. Proofs and programs share the
same language, and correctness is guaranteed by a small trusted kernel that
independently type-checks every proof term.

The \Mathlib library~\cite{mathlib2020} provides formalized mathematics
including abstract algebra, analysis, and topology. For our development,
the key components are:
\begin{itemize}
  \item \texttt{NormedRing}, \texttt{NormedAlgebra}: typeclasses for
    normed rings and algebras with submultiplicative norms.
  \item \texttt{NormedSpace.exp}: the exponential function
    $a \mapsto \sum_{k} a^k/k!$ in a complete normed algebra.
  \item \texttt{intervalIntegral}: the Bochner integral over $[a,b]$,
    with the fundamental theorem of calculus.
  \item \texttt{HasDerivAt}: the Fr\'{e}chet derivative, with product rules
    and chain rules for normed algebra-valued functions.
\end{itemize}

\paragraph{Related formalization work.}
Formal verification has recently gained traction in mathematical physics and
quantum information.
Zhao and Yu~\cite{zhao2025chsh} formalized robust CHSH rigidity in \Lean,
discovering a gap in the original proof of McKague~(2012) during the
formalization process.
Meiburg et al.~\cite{meiburg2025steinlemma} formalized a generalized
quantum Stein's lemma, and
Tooby-Smith~\cite{toobysmith2026higgs} formalized the stability of the
two-Higgs-doublet potential, also identifying a literature error.
These projects share the pattern of using formal verification not merely to
certify known results but to expose subtle gaps in published proofs.
Our work extends this program to quantum simulation, providing the first
machine-checked proof of any product formula for operator exponentials.

\section{Proof Architecture}
\label{sec:proof}

Our formalization is organized into five independent tracks that compose to yield
the main theorem and its extensions. \Cref{fig:dag} shows the dependency
structure. We present the mathematical arguments here; the Lean formalization
is discussed in \cref{sec:formalization}.

\begin{figure}[ht]
\centering
\small
\[
\underbrace{
  \begin{array}{c}
  \texttt{Telescoping} \\[2pt]
  X^n - Y^n = \sum X^k(X{-}Y)Y^{n{-}1{-}k}
  \end{array}
}_{\text{Track 1}}
\quad
\underbrace{
  \begin{array}{c}
  \texttt{ExpBounds} \\[2pt]
  \|e^a\| \le e^{\|a\|}
  \end{array}
}_{\text{Track 2}}
\quad
\underbrace{
  \begin{array}{c}
  \texttt{ExpDivPow} \\[2pt]
  (e^{a/n})^n = e^a
  \end{array}
}_{\text{Track 4}}
\]
\[
\downarrow \qquad\qquad\qquad\qquad \downarrow
\]
\[
\underbrace{
  \begin{array}{c}
  \texttt{StepError} \\[2pt]
  \|e^a e^b - e^{a+b}\| \le 2\|a\|\|b\|\,e^{\|a\|+\|b\|}
  \end{array}
}_{\text{Track 3 (depends on Track 2)}}
\]
\[
\downarrow
\]
\[
\boxed{
  \begin{array}{c}
  \texttt{Assembly} \\[2pt]
  (e^{A/n}\,e^{B/n})^n \to e^{A+B} \quad [O(1/n)]
  \end{array}
}
\]
\[
\swarrow \qquad\qquad\qquad \searrow
\]
\[
\underbrace{
  \begin{array}{c}
  \texttt{StrangSplitting} \\[2pt]
  O(1/n^2)
  \end{array}
}_{\text{Extensions}}
\quad
\underbrace{
  \begin{array}{c}
  \texttt{Multi\{Operator,Strang\}} \\[2pt]
  m\text{-operator}
  \end{array}
}_{\text{Extensions}}
\quad
\underbrace{
  \begin{array}{c}
  \texttt{CommutatorScaling} \\[2pt]
  \|[B,A]\|\text{-scaling}
  \end{array}
}_{\text{Track 5 (independent)}}
\]
\caption{Dependency structure of the formalization.
  Tracks~1--4 compose into the main theorem; extensions branch from the
  assembly or develop independently.}
\label{fig:dag}
\end{figure}

\subsection{Modular decomposition}
\label{sec:modular}

The proof of the Lie--Trotter formula decomposes into four independent
components:
\begin{description}[leftmargin=1.5em]
  \item[Track~1 (Algebraic):] The telescoping identity and its norm bound.
  \item[Track~2 (Analytic):] Exponential series remainder estimates.
  \item[Track~3 (Step error):] The single-step approximation bound
    $\norm{e^a\,e^b - e^{a+b}} \le 2\norm{a}\norm{b}\,e^{\norm{a}+\norm{b}}$.
  \item[Track~4 (Connecting):] The identity $(e^{a/n})^n = e^a$.
\end{description}
Tracks 1--4 are combined in an assembly step to yield the convergence rate,
from which the limit theorem follows by a standard $\varepsilon$-$N$ argument.
The commutator-scaling bounds (Track~5) are an independent extension that
refines the step error.

\subsection{Telescoping and norm bounds}
\label{sec:telescoping}

The algebraic identity
\begin{equation}\label{eq:telescoping}
  X^n - Y^n = \sum_{k=0}^{n-1} X^k\,(X-Y)\,Y^{n-1-k}
\end{equation}
is proved by induction on $n$. The key step is factoring $Y$ out of the inner
sum to match the inductive hypothesis: writing $Y^{m-k} = Y^{m-1-k} \cdot Y$
converts $\sum_{k<m} X^k(X-Y)Y^{m-k}$ into $(\sum_{k<m} X^k(X-Y)Y^{m-1-k}) \cdot Y$.

The norm bound follows by the triangle inequality applied to each summand:
\begin{equation}\label{eq:norm-telescoping}
  \norm{X^n - Y^n} \le n\,\norm{X-Y}\,M^{n-1},
  \qquad M = \max\{\norm{X}, \norm{Y}\}.
\end{equation}

\subsection{Exponential series estimates}
\label{sec:exp-bounds}

We establish four bounds on the exponential in a complete normed algebra:
\begin{align}
  \norm{e^a} &\le e^{\norm{a}},
    \label{eq:norm-exp} \\
  \norm{e^a - 1} &\le e^{\norm{a}} - 1,
    \label{eq:norm-exp-sub-one} \\
  e^x - 1 - x &\le \tfrac{x^2}{2}\,e^x
    \quad (x \ge 0),
    \label{eq:exp-quadratic} \\
  \norm{e^a - 1 - a} &\le \tfrac{\norm{a}^2}{2}\,e^{\norm{a}}.
    \label{eq:norm-exp-quadratic}
\end{align}

The proofs use the power series representation $e^a = \sum_{n=0}^{\infty} a^n/n!$.
Bounds \eqref{eq:norm-exp}--\eqref{eq:norm-exp-sub-one} follow from
$\norm{\sum a^n/n!} \le \sum \norm{a}^n/n! = e^{\norm{a}}$ with appropriate
index shifts. Bound \eqref{eq:exp-quadratic} uses the termwise estimate
$x^{n+2}/(n+2)! \le (x^2/2) \cdot x^n/n!$ via the factorial inequality
$2 \cdot n! \le (n+2)!$. Bound \eqref{eq:norm-exp-quadratic} combines
\eqref{eq:exp-quadratic} with the same tsum comparison technique as
\eqref{eq:norm-exp}.

\begin{remark}
  Bound \eqref{eq:norm-exp} was not available in \Mathlib for general
  Banach algebras; \Mathlib only had the specialized result
  \texttt{Complex.norm\_exp\_le\_exp\_norm} for $\mathbb{C}$.
  Our proof from the power series fills this gap.
\end{remark}

\subsection{Step error via algebraic factorization}
\label{sec:step-error}

The single-step bound is the most delicate component. Rather than expanding
$e^a\,e^b$ and $e^{a+b}$ to second order and bounding cross-terms directly
(which is tractable on paper but unwieldy in a proof assistant due to
non-commutative bookkeeping), we use an algebraic factorization:
\begin{equation}\label{eq:factorization}
  e^a\,e^b - e^{a+b}
    = (e^a - 1)(e^b - 1) - \bigl(e^{a+b} - e^a - e^b + 1\bigr).
\end{equation}
The first term is bounded by $(e^{\norm{a}}-1)(e^{\norm{b}}-1)$ via
\eqref{eq:norm-exp-sub-one}. For the second (cross) term, we prove by
induction on the power $m$ that
\begin{equation}\label{eq:cross-term}
  \norm{(a+b)^m - a^m - b^m}
    \le (\norm{a}+\norm{b})^m - \norm{a}^m - \norm{b}^m
    \qquad (m \ge 1),
\end{equation}
using the non-commutative recurrence
$(a+b)^{m+1} - a^{m+1} - b^{m+1} = (a+b)\bigl((a+b)^m - a^m - b^m\bigr) + ab^m + ba^m$.
Summing over the exponential series, the cross term is also bounded by
$(e^{\norm{a}}-1)(e^{\norm{b}}-1)$, so by the triangle inequality:
\begin{equation}\label{eq:step-bound}
  \norm{e^a\,e^b - e^{a+b}}
    \le 2(e^{\norm{a}}-1)(e^{\norm{b}}-1)
    \le 2\norm{a}\norm{b}\,e^{\norm{a}+\norm{b}},
\end{equation}
where the last step uses $(e^s-1)(e^t-1) \le st\,e^{s+t}$ for $s,t \ge 0$.

\begin{remark}
  The factorization \eqref{eq:factorization} was discovered during
  formalization: the standard second-order expansion approach required
  tracking non-commutative multinomial cross-terms, which was intractable
  in Lean. The algebraic factorization reduces the problem to two
  independent bounds, each handled by series comparison.
\end{remark}

\subsection{Assembly and convergence}
\label{sec:assembly}

Setting $P = e^{A/n}\,e^{B/n}$ and $Q = e^{(A+B)/n}$:
\begin{enumerate}
  \item $Q^n = e^{A+B}$ by the identity $(e^{a/n})^n = e^a$
    (\leanref{ExpDivPow.lean}{31}{exp\_div\_pow}).
  \item $\norm{P-Q} \le 2\norm{A}\norm{B}/n^2 \cdot e^{(\norm{A}+\norm{B})/n}$
    by the step error with $a = A/n$, $b = B/n$.
  \item $\max\{\norm{P},\norm{Q}\} \le e^{(\norm{A}+\norm{B})/n}$ by
    \eqref{eq:norm-exp}.
  \item Telescoping gives
    $\norm{P^n - Q^n} \le n \cdot \norm{P-Q} \cdot M^{n-1}$.
  \item The factors $e^{(\norm{A}+\norm{B})/n}$ combine:
    $e^{s/n} \cdot (e^{s/n})^{n-1} = (e^{s/n})^n = e^s$, yielding
    \[
      \norm{(e^{A/n}\,e^{B/n})^n - e^{A+B}}
        \le \frac{2\norm{A}\norm{B}\,e^{\norm{A}+\norm{B}}}{n}.
    \]
\end{enumerate}
The main theorem $\lim_{n\to\infty} (e^{A/n}\,e^{B/n})^n = e^{A+B}$ follows
by a standard $\varepsilon$-$N$ argument using the Archimedean property.

\subsection{Duhamel approach to commutator scaling}
\label{sec:commutator}

The commutator-scaling bound \eqref{eq:comm-first} replaces the product
$\norm{A}\norm{B}$ with $\norm{\comm{B}{A}}$. Our proof uses a
variation-of-parameters (Duhamel) technique based on the fundamental theorem
of calculus, avoiding ODE uniqueness theory entirely.

\paragraph{Integral error representation.}
Define the ``conjugated Trotter product''
\begin{equation}\label{eq:w-def}
  w(\tau) = e^{-\tau(A+B)}\,e^{\tau B}\,e^{\tau A}.
\end{equation}
By the product rule for Fr\'{e}chet derivatives in a normed algebra,
\begin{equation}\label{eq:w-deriv}
  w'(\tau) = e^{-\tau(A+B)}\,\bigl(e^{\tau B}\,A - A\,e^{\tau B}\bigr)\,e^{\tau A}
  = e^{-\tau(A+B)}\,\comm{e^{\tau B}}{A}\,e^{\tau A}.
\end{equation}
The key cancellation is that the ``free'' terms involving $A+B$, $B$, and $A$
sum to zero: $-(A+B) + B + A = 0$. The FTC-2 gives
\begin{equation}\label{eq:integral-error}
  e^{tB}\,e^{tA} - e^{t(A+B)}
    = \int_0^t e^{(t-\tau)(A+B)}\,\comm{e^{\tau B}}{A}\,e^{\tau A}\,d\tau,
\end{equation}
after left-multiplying by $e^{t(A+B)}$ (which commutes with itself).

\paragraph{Extracting the commutator.}
The integrand contains $\comm{e^{\tau B}}{A} = e^{\tau B}A - Ae^{\tau B}$,
which involves the \emph{exponential} of $B$. To extract the bare commutator
$\comm{B}{A}$, we apply FTC-2 a second time to the conjugation map
$s \mapsto e^{sB}\,A\,e^{-sB}$:
\begin{equation}\label{eq:conj-ftc}
  e^{\tau B}\,A\,e^{-\tau B} - A
    = \int_0^{\tau} e^{sB}\,\comm{B}{A}\,e^{-sB}\,ds.
\end{equation}
This gives $\norm{e^{\tau B}\,A\,e^{-\tau B} - A}
  \le \norm{\comm{B}{A}}\,\abs{\tau}\,e^{2\abs{\tau}\norm{B}}$
by bounding the integrand and using
$\norm{\int_0^\tau f} \le C\abs{\tau}$.
Right-multiplying by $e^{\tau B}$ yields
$\norm{\comm{e^{\tau B}}{A}} \le \norm{\comm{B}{A}}\,\abs{\tau}\,e^{3\abs{\tau}\norm{B}}$.

\paragraph{Final bound.}
Substituting into \eqref{eq:integral-error} and bounding the integral:
\[
  \norm{e^{tB}\,e^{tA} - e^{t(A+B)}}
    \le \int_0^t e^{(t-\tau)(\norm{A}+\norm{B})} \cdot
      \norm{\comm{B}{A}}\,\tau\,e^{3\tau\norm{B}} \cdot e^{\tau\norm{A}}\,d\tau.
\]
Simplifying the exponential factors and evaluating $\int_0^t \tau\,d\tau = t^2/2$
gives the claimed bound \eqref{eq:comm-first}.

\begin{remark}
  The Duhamel approach requires only the product rule for Fr\'{e}chet
  derivatives, the FTC-2 for Bochner integrals, and basic norm estimates---all
  available in \Mathlib. The alternative approach via Gronwall's inequality
  would require ODE existence/uniqueness theory, adding approximately 40 lines
  of additional Lean infrastructure.
\end{remark}

\subsection{Second-order Strang commutator scaling}
\label{sec:strang}

For the Strang splitting $S_2(t) = e^{tA/2}\,e^{tB}\,e^{tA/2}$, we prove
\begin{equation}\label{eq:strang-bound}
  \norm{S_2(t) - e^{tH}}
    \le \frac{\norm{\comm{B}{\comm{B}{A}}}}{12}\,t^3
      + \frac{\norm{\comm{A}{\comm{A}{B}}}}{24}\,t^3,
\end{equation}
where $H = A + B$ and we additionally require that $A$ and $B$ are
\emph{skew-adjoint} ($a^* = -a$), so that $\norm{e^{ta}} = 1$.

The proof applies the same Duhamel strategy to
$w(\tau) = e^{-\tau H}\,S_2(\tau)$, but now requires:
\begin{enumerate}
  \item A \emph{4-factor product rule} for $S_2(\tau) = e^{\tau A/2}\,e^{\tau B}\,e^{\tau A/2}$,
    yielding $w'(\tau) = e^{-\tau H}\,\mathcal{T}_2(\tau)\,S_2(\tau)$, where $\mathcal{T}_2$
    involves the commutators $\comm{A/2}{e^{\tau B}}$ and $\comm{e^{\tau B}}{A/2}$.

  \item \emph{First-order cancellation}: the leading commutator terms cancel
    exactly: $\comm{A/2}{B} + \comm{B}{A/2} = 0$.

  \item A \emph{double FTC} on $s \mapsto e^{sB}\,A\,e^{-sB}$, expanding to
    second order to extract $\comm{B}{\comm{B}{A}}$.

  \item The \emph{``subtract-constant-at-$\tau$'' trick}: to bound
    $\int_0^\tau F(s)\,ds - \tau\,F(\tau)$, rewrite as
    $\int_0^\tau (F(s) - F(\tau))\,ds$ and bound $\norm{F(s) - F(\tau)}
    \le C_A(\tau - s)$.
    This avoids Fubini's theorem or integration by parts.
\end{enumerate}

\paragraph{The 4-factor product rule.}
Write $A' = A/2$ and $H = A + B$. The conjugated Strang product
\[
  w(\tau) = e^{-\tau H}\,e^{\tau A'}\,e^{\tau B}\,e^{\tau A'}
\]
has four exponential factors, each contributing a derivative term via the
Leibniz rule. Using $d/d\tau[e^{\tau X}] = X\,e^{\tau X}$ and collecting:
\begin{equation}\label{eq:4factor}
  w'(\tau) = e^{-\tau H}\,\mathcal{T}_2(\tau)\,S_2(\tau),
\end{equation}
where the \emph{Strang residual} is
\[
  \mathcal{T}_2(\tau)
    = \bigl(e^{\tau A'}\,B\,e^{-\tau A'} - B\bigr)
      + e^{\tau A'}\bigl(e^{\tau B}\,A'\,e^{-\tau B} - A'\bigr)e^{-\tau A'}.
\]
The key algebraic step is verifying that the ``free'' derivative terms
$-H + A' + B + A'$ sum to zero (since $A' + A' = A$ and $H = A + B$).
In Lean, the sum of the four derivative contributions is factored as
$-(e^{-\tau H}) \cdot (H - A' - A' - B) \cdot S_2(\tau)$, where
$H - A' - A' - B = 0$ is discharged by \texttt{abel}. The free-ring
factoring is handled by \texttt{noncomm\_ring}.

\paragraph{Coefficient conversion.}
Since $A' = A/2$, the double commutators in the final bound involve
half-scaled operators:
$\comm{B}{\comm{B}{A'}} = \frac{1}{2}\comm{B}{\comm{B}{A}}$ and
$\comm{A'}{\comm{A'}{B}} = \frac{1}{4}\comm{A}{\comm{A}{B}}$.
The scalar factors are extracted through $\comm{cX}{Y} = c\,\comm{X}{Y}$
(valid for central scalars) using \texttt{Algebra.smul\_mul\_assoc} and
\texttt{Algebra.mul\_smul\_comm}, then \texttt{norm\_smul}.
Combined with the integration constants ($1/6$ from the cubic integral
and $1/2$ from the double FTC), this produces the final coefficients
$1/12$ and $1/24$ in \eqref{eq:strang-bound}.

\section{Formalization in Lean~4}
\label{sec:formalization}

\subsection{Project structure and statistics}
\label{sec:project-structure}

The formalization is organized into 41~Lean source files, grouped by track
in \cref{tab:files}. Each file is self-contained with explicit imports,
enabling parallel development and independent verification.

\begin{table}[ht]
\centering
\footnotesize
\setlength{\tabcolsep}{4pt}
\begin{tabular}{@{}lrl@{}}
\toprule
\textbf{File} & \textbf{Lines} & \textbf{Content} \\
\midrule
\multicolumn{3}{@{}l}{\textit{Lie--Trotter core (Tracks 1--4, \cref{sec:proof})}} \\
\leansrc{Telescoping.lean}       & 106  & Algebraic telescoping + norm bound \\
\leansrc{ExpBounds.lean}         & 212  & Exponential series remainder estimates \\
\leansrc{StepError.lean}         & 674  & Quadratic step error (algebraic factorization) \\
\leansrc{ExpDivPow.lean}         & 54   & $(e^{a/n})^n = e^a$ \\
\leansrc{Assembly.lean}          & 146  & $O(1/n)$ convergence + main theorem \\
\midrule
\multicolumn{3}{@{}l}{\textit{Strang, multi-operator, commutator scaling (Track~5)}} \\
\leansrc{StrangSplitting.lean}   & 492  & Strang splitting, $O(1/n^2)$ \\
\leansrc{MultiOperator.lean}     & 288  & Multi-operator first-order \\
\leansrc{MultiStrang.lean}       & 428  & Multi-operator Strang \\
\leansrc{CommutatorScaling.lean} & 379  & Duhamel commutator-scaling \\
\leansrc{MultiCommutatorScaling.lean}       & 128  & Multi-operator comm.\ scaling \\
\leansrc{StrangCommutatorScaling.lean}      & 733  & Strang double-commutator scaling \\
\leansrc{MultiStrangCommutatorScaling.lean} & 174  & Multi-op.\ Strang comm.\ scaling \\
\leansrc{HigherCommutator.lean}            & 243  & Triple-FTC: extracts $[B,[B,[B,A]]]$ \\
\leansrc{StrangCommutatorScalingTight.lean} & 603 & Tighter Strang via norm-of-difference $D$; $\norm{D}$ domination \\
\leansrc{StrangTotalErrorCommScaling.lean}  & 184 & Commutator-scaled Strang \emph{total} error \\
\midrule
\multicolumn{3}{@{}l}{\textit{Suzuki $S_4$ track (\cref{sec:suzuki4})}} \\
\leansrc{Suzuki4.lean}                    & 466  & $S_4$ definition; generic $O(1/n^2)$ rate \\
\leansrc{Suzuki4FullDuhamel.lean}         & 445  & $S_4$ $O(t^3)$ via 5-$S_2$ telescoping \\
\leansrc{Suzuki4CommutatorScaling.lean}   & 90   & \texttt{suzuki4Exp} definition \\
\leansrc{Suzuki4Deriv.lean}               & 96   & Derivative scaffolding \\
\leansrc{Suzuki4HasDerivAt.lean}          & 136  & Module~1: \texttt{HasDerivAt} for 12-factor $w_4$ \\
\leansrc{Suzuki4Module2.lean}             & 167  & Module~2: FTC-2 bridge \\
\leansrc{Suzuki4Module3.lean}             & 184  & Module~3: FTC-2 reduction \\
\leansrc{Suzuki4Module4.lean}             & 175  & Module~4a: continuity of $w_4'$ \\
\leansrc{Suzuki4DerivExplicit.lean}       & 979  & Module~4b: explicit derivative; Phase 1--3 \\
\leansrc{Suzuki4StrangBlocks.lean}        & 120  & $S_4$ = 5 Strang blocks + cubic cancellation \\
\leansrc{Suzuki4MultinomialExpand.lean}   & 787  & Multinomial formulas; h2 unconditional, h3 under Suzuki cubic \\
\leansrc{Suzuki4Phase5.lean}              & 782  & Taylor reduction + Leibniz bridges + CAPSTONE \\
\leansrc{Suzuki4OrderFive.lean}           & 440  & $O(t^5)$ abstract-form target \\
\leansrc{Suzuki4ChildsForm.lean}          & 216  & Childs Prop pf4\_bound\_2term framework \\
\leansrc{Suzuki4BchBound.lean}            & 495  & SLICE~1: single-step $O(\abs{\tau}^5)$ bound \\
\leansrc{TaylorMatch.lean}                & 331  & SLICE~2: generic Taylor-match-from-norm \\
\leansrc{Suzuki4ViaBCH.lean}              & 1\,436 & SLICE~3 wiring + L1/L2/L3/L4 BCH-derived bounds \\
\leansrc{Suzuki4Convergence.lean}         & 398  & Total error $O(1/n^4)$ + convergence \\
\leansrc{Suzuki4TightConvergence.lean}    & 169  & Commutator-scaled $O(1/n^4)$ \\
\leansrc{Suzuki4UnitaryTotalError.lean}   & 274  & Skew-adjoint total error, no growth factor \\
\leansrc{Suzuki4Commute.lean}             & 228  & Commuting case: $S_4$ exact, error $\equiv 0$ \\
\leansrc{Suzuki4GapClosers.lean}          & 205  & General-$t$ step; imaginary-time corollaries \\
\leansrc{TrotterStepCount.lean}           & 392  & $\varepsilon$-form step counts; explicit per-$n$ bounds \\
\leansrc{MatrixCorollaries.lean}          & 157  & Matrix specializations (spectral norm) \\
\leansrc{PrefactorStrict.lean}            & 111  & Strict $\gamma_i < \alpha_i$; gap $\gamma_i \le \alpha_i/8$ \\
\midrule
\multicolumn{3}{@{}l}{\textit{Root}} \\
\href{\repourl/LieTrotter.lean}{\texttt{LieTrotter.lean}} & 73 & Import aggregator \\
\midrule
\textbf{Total} & \textbf{14\,196} & \textbf{0 \texttt{sorry} axioms,
0 own theorem-level BCH-interface axioms} \\
\bottomrule
\end{tabular}
\caption{Source files in the formalization, grouped by track.
\texttt{Suzuki4.lean} carries the generic $O(1/n^2)$ rate, which holds for
every $n \ge 1$ and needs no cubic condition;
\texttt{Suzuki4Convergence.lean} upgrades the \emph{same} object to $O(1/n^4)$
under the Suzuki cubic condition (\cref{sec:s4-convergence}), and
\texttt{Suzuki4TightConvergence.lean} does so with a commutator-scaled leading
coefficient.}
\label{tab:files}
\end{table}

The development builds against Lean~4.29.0-rc8 and current \Mathlib.
The companion
\href{https://github.com/Jue-Xu/Lean-BCH}{Lean-BCH} project is imported
as a git dependency, supplying the symmetric three-factor BCH identities,
the two-factor palindromic quintic remainder, and the five-factor
palindromic quintic identification used in the Suzuki~$S_4$ track.
Lean-BCH is sorry-free as of revision~\texttt{d455ff0} (2026-05-19),
which closed the previous B1.c (\texttt{symmetric\_bch\_quintic\_sub\_poly\_axiom})
dependency; its single-step $O(\abs{\tau}^5)$ payoff lemma anchors
our SLICE-1 bound in \leansrc{Suzuki4BchBound.lean}. Two private
septic axioms remain in Lean-BCH and gate only the optional Level~4
local small-$\tau$ ($\tau^7$) refinement (see \cref{sec:caveats}).
Build time for the full project is under 2~minutes on a standard
laptop after \Mathlib oleans are cached (\texttt{lake exe cache get}).

\subsection{Type-level setup}
\label{sec:type-setup}

Our theorems are stated for a general complete normed algebra $\A$ over a
field $\mathbb{K} \in \{\mathbb{R}, \mathbb{C}\}$:

\begin{leancode}
variable {𝕂 : Type*} [RCLike 𝕂]
variable {𝔸 : Type*} [NormedRing 𝔸] [NormedAlgebra 𝕂 𝔸]
  [NormOneClass 𝔸] [CompleteSpace 𝔸]
\end{leancode}

The typeclass \texttt{NormOneClass} requires $\norm{1} = 1$, which is needed
for \texttt{norm\_pow\_le} in current \Mathlib. The main theorem is
(\leanref{Assembly.lean}{121}{lie\_trotter}):

\begin{leancode}
theorem lie_trotter (A B : 𝔸) :
    Filter.Tendsto
      (fun n : ℕ => (exp ((n : 𝕂)⁻¹ • A) *
                      exp ((n : 𝕂)⁻¹ • B)) ^ n)
      atTop (nhds (exp (A + B)))
\end{leancode}

Note the scalar action $\bullet$: the step is $n^{-1}$ \emph{acting on} $A$
in the $\mathbb{K}$-module $\A$, not a ring product with a coerced $n^{-1} \in
\A$. The two coincide in a unital algebra but are distinct terms in Lean, and
the derivative and integral lemmas we rely on are all stated for $\bullet$.

\paragraph{The \texttt{include} pattern.}
Since \texttt{NormedSpace.exp} no longer takes a field parameter in recent
\Mathlib, the variable $\mathbb{K}$ does not appear in the type of theorems
involving \texttt{exp}---in \texttt{norm\_exp\_le} below, both sides mention
only $\A$ and $\mathbb{R}$. Lean~4 drops unused type-level variables, so we
must write \texttt{include $\mathbb{K}$ in} before each lemma whose proof (but
not statement) requires $\mathbb{K}$
(e.g., \leanref{ExpBounds.lean}{72}{norm\_exp\_le}, whose proof rewrites with
\texttt{exp\_tsum\_form ($\mathbb{K}$ := $\mathbb{K}$)}):

\begin{leancode}
include 𝕂 in
lemma norm_exp_le (a : 𝔸) : ‖exp a‖ ≤ Real.exp ‖a‖ := by ...
\end{leancode}

\subsection{Telescoping in Lean}
\label{sec:lean-telescoping}

The telescoping identity~\eqref{eq:telescoping} is proved by induction on $n$
using \texttt{Finset.sum\_range\_succ} to peel off the last term
(\leanref{Telescoping.lean}{32}{telescoping\_direct}):

\begin{leancode}
theorem telescoping_direct (X Y : R) (n : ℕ) :
    (∑ k ∈ Finset.range n,
      X ^ k * (X - Y) * Y ^ (n - 1 - k)) =
      X ^ n - Y ^ n := by
  induction n with
  | zero => simp
  | succ m ih =>
    rw [Finset.sum_range_succ]
    -- Factor Y from the inner sum to get the
    -- inductive hypothesis
    ...
    noncomm_ring [pow_succ]
\end{leancode}

Here \texttt{R} is any \texttt{Ring}: the telescoping identity is purely
algebraic, so this file assumes no norm and no completeness.

The final step uses \texttt{noncomm\_ring}, which handles polynomial identities
in non-commutative rings. The standard \texttt{ring} tactic silently fails on
non-commutative goals---it produces an unsolved subgoal rather than an error,
which cost us significant debugging time early in the project.

\subsection{Exponential bounds: filling Mathlib gaps}
\label{sec:lean-exp-bounds}

The bound $\norm{e^a} \le e^{\norm{a}}$ was not available in \Mathlib for
general Banach algebras. Our proof works from the tsum representation
(\leanref{ExpBounds.lean}{72}{norm\_exp\_le}):

\begin{leancode}
include 𝕂 in
lemma norm_exp_le (a : 𝔸) : ‖exp a‖ ≤ Real.exp ‖a‖ := by
  rw [exp_tsum_form (𝕂 := 𝕂), real_exp_eq_tsum]
  calc ‖∑' n, (((Nat.factorial n : 𝕂)⁻¹) • a ^ n : 𝔸)‖
      ≤ ∑' n, ‖...‖         -- norm_tsum_le_tsum_norm
    _ ≤ ∑' n, ‖a‖ ^ n / n ! -- norm_exp_term_le
    _ = Real.exp ‖a‖        -- Real.hasSum_exp
\end{leancode}

The proof chains three standard ingredients: pulling the norm inside the tsum,
bounding each term $\norm{a^n/n!} \le \norm{a}^n/n!$ via submultiplicativity,
and recognizing the resulting series as $e^{\norm{a}}$.

\subsection{Step error: algebraic factorization in Lean}
\label{sec:lean-step-error}

The factorization~\eqref{eq:factorization} is verified by \texttt{noncomm\_ring}.
The cross-term bound~\eqref{eq:cross-term} is proved by induction using the
non-commutative recurrence
(\leanref{StepError.lean}{43}{norm\_pow\_add\_sub\_pow\_sub\_pow}):

\begin{leancode}
private lemma norm_pow_add_sub_pow_sub_pow (a b : 𝔸) :
    ∀ m : ℕ, 1 ≤ m →
      ‖(a + b) ^ m - a ^ m - b ^ m‖ ≤
        (‖a‖ + ‖b‖) ^ m - ‖a‖ ^ m - ‖b‖ ^ m := by
  intro m hm
  induction m with
  | zero => omega
  | succ n ih => ...
\end{leancode}

At 674~lines, \leansrc{StepError.lean} is the longest file in the basic
Lie--Trotter development (excluding commutator-scaling extensions).
The inductive cross-term bound accounts for roughly half of this length,
because each step requires explicit \texttt{norm\_mul\_le} applications and
\texttt{gcongr} for the inductive comparison.

\subsection{Commutator scaling: calculus in Lean}
\label{sec:lean-commutator}

The commutator-scaling proofs are the most technically demanding part of the
formalization, as they require doing analysis---derivatives, interval integrals,
FTC-2---in a non-commutative Banach algebra.

\paragraph{Derivative computations.}
\Mathlib's \texttt{hasDerivAt\_exp\_smul\_const'} supplies the
derivative $d/d\tau[e^{\tau B}] = B\,e^{\tau B}$, and the
product rule \texttt{HasDerivAt.mul} extends to normed algebras.
One subtlety: $(-\tau) \bullet B$ and $\tau \bullet (-B)$ are
definitionally distinct but equal; we must normalize before
applying derivative lemmas
(\leanref{CommutatorScaling.lean}{75}{hasDerivAt\_exp\_conj}):

\begin{leancode}
lemma hasDerivAt_exp_conj (A B : 𝔸) (s : ℝ) :
    HasDerivAt
      (fun u => exp (u • B) * A * exp ((-u) • B))
      (exp (s • B) * (B * A - A * B) *
       exp ((-s) • B)) s := by
  simp_rw [show ∀ u : ℝ, (-u) • B = u • (-B)
    from fun u => by rw [neg_smul, smul_neg]]
  ...
\end{leancode}

\paragraph{FTC-2 for integral representations.}
The conjugation identity~\eqref{eq:conj-ftc} and the integral error
representation~\eqref{eq:integral-error} are proved by
\texttt{integral\_eq\_sub\_of\_hasDerivAt}
(\leanref{CommutatorScaling.lean}{157}{exp\_conj\_sub\_eq\_integral}):

\begin{leancode}
theorem exp_conj_sub_eq_integral (A B : 𝔸) (τ : ℝ) :
    exp (τ • B) * A * exp ((-τ) • B) - A =
      ∫ s in (0 : ℝ)..τ,
        exp (s • B) * (B * A - A * B) *
        exp ((-s) • B) := by
  have hderiv : ∀ u ∈ Set.uIcc 0 τ,
      HasDerivAt ... := fun u _ =>
    hasDerivAt_exp_conj A B u
  have hint : IntervalIntegrable ... :=
    (continuous_exp_conj_deriv A B).intervalIntegrable 0 τ
  exact (integral_eq_sub_of_hasDerivAt hderiv hint).symm
\end{leancode}

\paragraph{Pulling constants through integrals.}
In a \texttt{NormedRing}, left-multiplication by a fixed element $c$ is
\emph{not} a scalar multiplication---it is ring multiplication. To derive
$c \cdot \int f = \int c \cdot f$, we use
\texttt{ContinuousLinearMap.intervalIntegral\_comp\_comm} with
$L = \texttt{ContinuousLinearMap.mul}\ \mathbb{R}\ \A\ c$, which views
left-multiplication as a continuous linear map.

\paragraph{Working over $\mathbb{R}$.}
The commutator-scaling files work over $\texttt{NormedAlgebra}\ \mathbb{R}\ \A$
directly, rather than a general $\mathbb{K}$. The reason is that the
\texttt{HasDerivAt}/\texttt{intervalIntegral} machinery requires
$\texttt{SMul}\ \mathbb{R}\ \A$, which is \emph{not} automatically synthesized
from $\texttt{[RCLike K] [NormedAlgebra K A]}$. For applications with
$\mathbb{K} = \mathbb{C}$, one can use
$\texttt{NormedAlgebra.restrictScalars}\ \mathbb{R}\ \mathbb{C}\ \A$.

\subsection{Multi-operator and Suzuki extensions}
\label{sec:lean-extensions}

The multi-operator formulas are proved by induction on the list of operators.
The palindromic product is defined recursively
(\leanref{MultiStrangCommutatorScaling.lean}{42}{palindromicProd}):

\begin{leancode}
def palindromicProd : List 𝔸 → 𝔸
  | [] => 1
  | [a] => exp a
  | a :: rest =>
    exp ((1/2 : ℝ) • a) *
    palindromicProd rest *
    exp ((1/2 : ℝ) • a)
\end{leancode}

The step size is not a parameter: it is folded into the operators by the
caller, which passes the list $[t \bullet A_1, \ldots, t \bullet A_m]$. This
keeps the recursion on a single argument, so the induction is on the list
alone. The inductive step peels off the outermost $e^{A_1/2}$ factors and
applies the two-operator Strang bound to the pair
$(A_1, A_2 + \cdots + A_m)$, then recurses on the inner product.

\section{Discussion}
\label{sec:discussion}

\subsection{AI assistance in the development process}
\label{sec:ai}

We used LLM-based coding assistants (Claude Code, Anthropic) to accelerate
formalization, in the same general direction as recent work on agentic
autoformalization~\cite{ren2026merlean}.

Our workflow followed a \emph{sorry-driven development} pattern:
\begin{enumerate}
  \item State all theorem signatures with \texttt{sorry} placeholders.
  \item Verify that the signatures compose correctly (the main theorem
    type-checks with all intermediates as \texttt{sorry}).
  \item Fill \texttt{sorry}s bottom-up, from leaves of the dependency DAG
    to the root.
\end{enumerate}
This approach worked well with LLM agents: independent leaves
(e.g., the four exponential bounds B1--B4) could be assigned to parallel
agents, each given the exact signature to fill and the list of available
lemmas.

\textbf{Where agents excelled.} ``Fill this \texttt{sorry} given these
lemmas''---tasks with a clear specification and known
ingredients---compiled on the first attempt when delegated to agents.
Most of the exponential bounds and the assembly step were completed this way.

\textbf{Where agents struggled.} ``Find the right proof approach''---tasks
requiring mathematical insight, such as the Strang commutator-scaling argument
or the algebraic factorization for the step error---required human-driven
exploration. The agent's tendency to attempt minor variations of a failing
approach (rather than reconsidering the overall strategy) became a bottleneck
for these tasks.

\textbf{A cautionary pattern.} As noted by Zhao et al.~\cite{zhao2025chsh},
LLM agents sometimes quietly strengthen hypotheses to make a proof trivial.
The \texttt{sorry}-first workflow guards against this: since all signatures are
fixed before proof filling begins, adding hypotheses breaks the type-checking
of downstream theorems.

\subsection{Mathlib ecosystem}
\label{sec:mathlib}

Several gaps in \Mathlib were encountered and filled during this project:
\begin{itemize}
  \item $\norm{e^a} \le e^{\norm{a}}$ for general Banach algebras
    (\leanref{ExpBounds.lean}{72}{norm\_exp\_le}).
  \item $e^x - 1 - x \le x^2/2 \cdot e^x$ for $x \ge 0$
    (\leanref{ExpBounds.lean}{129}{exp\_sub\_one\_sub\_bound\_real}).
  \item $\norm{e^a - 1} \le e^{\norm{a}} - 1$
    (\leanref{ExpBounds.lean}{88}{norm\_exp\_sub\_one\_le}).
\end{itemize}
These are natural candidates for upstreaming to \Mathlib.

Working over a general field $\mathbb{K}$ (via \texttt{RCLike}) introduced
friction with the calculus infrastructure, which assumes $\mathbb{R}$-module
structure. Our pragmatic solution---working over $\mathbb{R}$ directly in the
commutator-scaling files---is effective but means users must apply
\texttt{restrictScalars} for $\mathbb{C}$ applications. A unified treatment
would require additional typeclass instances in \Mathlib.

\subsection{Formalization-driven proof design}
\label{sec:formalization-insights}

The formalization process influenced proof design in several ways:

\textbf{Algebraic factorization.}
The standard approach to the step error---expanding $e^a\,e^b$ to second order
and bounding cross-terms---requires tracking non-commutative multinomial
coefficients. In Lean, this involves dozens of \texttt{norm\_mul\_le} and
\texttt{gcongr} steps for each cross-term. The algebraic factorization
(\cref{sec:step-error}) splits the problem into two independent bounds,
each handled by series comparison, reducing the proof length by approximately
half. This factorization is arguably cleaner than the standard approach even
on paper.

\textbf{Duhamel without Gronwall.}
The commutator-scaling proof was designed with Lean feasibility in mind.
The standard approach uses Gronwall's inequality or ODE uniqueness, which
would require formalizing existence and uniqueness for linear ODEs in
Banach spaces---theory not yet available in \Mathlib. The Duhamel/FTC approach
uses only the product rule and FTC-2, both already available.

\textbf{\texttt{noncomm\_ring} vs.\ \texttt{ring}.}
The \texttt{ring} tactic in Lean silently fails on non-commutative goals:
it produces an unsolved subgoal rather than an error message. We discovered
this only after extended debugging. The \texttt{noncomm\_ring} tactic handles
free non-commutative ring identities correctly but cannot use commutativity
hypotheses (e.g., $[a, e^a] = 0$). Our workaround: use \texttt{set} to make
exponentials opaque, rewrite commutativity as equalities, then apply
\texttt{noncomm\_ring} on the resulting free-ring identity.

\subsection{Relation to quantum simulation}
\label{sec:quantum}

Our results are stated for general complete normed algebras, which include
the bounded operators $B(\mathcal{H})$ on a Hilbert space as a special case.
Specialization to finite-dimensional matrix algebras
$M_d(\mathbb{C})$ is straightforward via \Mathlib's
\texttt{Matrix.instNormedRing}.

The commutator-scaling bounds are physically motivated by Hamiltonian
simulation: for a lattice Hamiltonian $H = \sum_j h_j$ where each $h_j$ acts
locally, the commutator $\comm{h_j}{h_k}$ vanishes for distant terms.
Our formalized bounds could serve as certified building blocks for verified
resource estimation in quantum algorithms, once combined with gate-count
analysis for individual exponentials.

\section{Conclusion and Future Work}
\label{sec:conclusion}

We have presented, to our knowledge, the first kernel-checked proof, with no
project-specific axioms, of norm convergence of the Lie--Trotter sequence for
arbitrary elements of a complete normed algebra, together with higher-order
extensions and commutator-scaling error bounds, in Lean~4 with the \Mathlib
library.
The development comprises over 14\,000 lines with zero \texttt{sorry} axioms
and zero own theorem-level BCH-interface axioms. Levels~1--4 are fully
proved at the project level: at Lean-BCH revision \texttt{05e8c52},
\texttt{\#print axioms} returns only \texttt{[propext,
Classical.choice, Quot.sound]} on both the $\tau^5$ headlines and the
Level~4 local small-$\tau$ ($\tau^7$) refinement. The development covers:
\begin{itemize}
  \item First-order Lie--Trotter convergence at rate $O(1/n)$;
  \item Symmetric Strang splitting at rate $O(1/n^2)$;
  \item Multi-operator generalizations of both formulas;
  \item First- and second-order commutator-scaling bounds via a
    Duhamel/FTC technique;
  \item A second-order Strang bound via norm-of-difference that captures
    signed cancellations invisible to the sum-of-norms form, and whose
    leading coefficient is machine-checked never to exceed it---strictly
    smaller whenever the two double commutators partially cancel, a gain
    measured at $9$--$50\%$ across standard spin-chain benchmarks
    (\cref{apd:tighter});
  \item A fourth-order Suzuki $S_4$ track with four CAS-assisted
    BCH-based bounds, including a replacement of the 2021 prefactors of
    Childs et al.\ (rigorous, but not proven tight) by CAS-certified
    coefficients smaller term by term by 8.6$\times$--64$\times$---the
    comparison machine-checked at the leading coefficient, with full-bound
    domination on a neighbourhood of zero under a strict-gap
    hypothesis---and an existential-$\delta$ small-$\tau$ refinement via
    an explicit $\tau^7$ correction;
  \item Fourth-order Suzuki convergence at rate $O(1/n^4)$, obtained by
    compounding the single-step $O(t^5)$ bound over $n$ steps --- and, from
    the sharp step bound, with an explicit commutator-scaled leading
    coefficient.
\end{itemize}
The first, second, and fourth-order convergence theorems together form a
complete verified hierarchy $O(1/n) \to O(1/n^2) \to O(1/n^4)$: for each
integrator the development supplies both a single-step error bound and the
compounded statement that the $n$-step product converges to $e^{t(A+B)}$ at the
corresponding rate.

The Duhamel approach to commutator scaling---avoiding ODE uniqueness in favor
of the fundamental theorem of calculus---may be of independent interest as a
proof strategy for related product formula bounds.

\paragraph{Future directions.}
Several extensions are natural:
\begin{enumerate}
  \item \emph{Mathlib contributions.}
    The exponential norm bounds (\leanref{ExpBounds.lean}{72}{norm\_exp\_le},
    \leanref{ExpBounds.lean}{88}{norm\_exp\_sub\_one\_le},
    \leanref{ExpBounds.lean}{129}{exp\_sub\_one\_sub\_bound\_real})
    are general-purpose lemmas suitable for upstreaming to \Mathlib.

  \item \emph{Trotter-native and higher-order BCH derivations.}
    The three-factor symmetric BCH identities, the two-factor
    palindromic quintic remainder, and the five-factor palindromic
    quintic and septic identifications are formalized in the Lean-BCH
    companion project at revision~\texttt{05e8c52}. A Trotter-native
    operator-algebra proof of the order-4 Taylor-coefficient identity
    (via \texttt{scripts/compute\_sumQuadCorr\_factored.py}) remains a
    useful independent derivation, while extending the certified BCH
    expansion to order nine would sharpen the local remainder.

  \item \emph{Specialization to quantum systems.}
    Connecting our abstract bounds to concrete gate-count estimates for
    Hamiltonian simulation on quantum computers would bridge the gap between
    formal mathematics and verified quantum algorithms.

  \item \emph{Unbounded operators.}
    The original Trotter--Kato theorem applies to strongly continuous
    semigroups on Banach spaces, covering unbounded generators.
    Formalizing this extension would require substantial additional
    functional analysis infrastructure in \Mathlib.
\end{enumerate}

More broadly, our project demonstrates that substantial results in operator
analysis---involving non-commutative algebra, power series manipulation,
Fr\'{e}chet calculus, and Bochner integration---are within reach of current
proof assistant technology.
We hope this work encourages further formalization of the mathematical
foundations underlying quantum algorithms.


\section*{Acknowledgments}
The Lean formalization was developed with assistance from
Claude Code (Anthropic), which was used both as a coding assistant for
filling proof obligations and as a writing assistant in the preparation
of this manuscript.

\bibliographystyle{plainnat}
\bibliography{references,ref_aps}
\end{document}
